\documentclass{article}
\usepackage[utf8]{inputenc}
\usepackage{geometry}
\geometry{a4paper, margin=1in}

\title{Evaluación de Respuestas - Nathaniel}
\author{Evaluado por Prof. CESAR RODRIGUEZ}
\date{\today}

\begin{document}

\maketitle

\section*{Análisis de las Respuestas del Cuestionario}

A continuación se presenta una evaluación de las respuestas proporcionadas por Nathaniel.

\begin{enumerate}
    \item \textbf{Concepto (Función de registros MX y NS):}
    
    \textbf{Evaluación: Acertada.}
    
    \textit{Comentario:} La respuesta define correctamente la función de un registro MX (Mail 
    Exchange) como el encargado de especificar los servidores de correo para un dominio. 
    Sin embargo, la respuesta omite la definición del registro NS. Aunque la parte respondida es 
    correcta, la respuesta está incompleta.

    \item \textbf{Análisis de Código (Prioridad en registros MX):}
    
    \textbf{Evaluación: Acertada.}
    
    \textit{Comentario:} Aquí explicas de manera correcta que el número de preferencia en los 
    registros MX indica la prioridad de entrega. Describes acertadamente que los servidores de 
    correo utilizan este valor para decidir a qué servidor contactar primero, lo cual es 
    fundamental para la redundancia y el correcto enrutamiento del correo.

    \item \textbf{Diseño de Software (Modularidad de la función):}
    
    \textbf{Evaluación: Regular.}
    
    \textit{Comentario:} La idea de "añadir un parámetro extra" es un paso en la dirección 
    correcta, pero tu respuesta es demasiado vaga. Una solución más modular y robusta, como se 
    esperaba, sería modificar la función para que \textbf{retorne} los datos (por ejemplo, en un 
    diccionario o una lista) en lugar de imprimirlos. De esta forma, la lógica de obtención de 
    datos se separa de la lógica de presentación, permitiendo que el código que llama a la función 
    decida si los imprime, los guarda en un archivo JSON, los formatea en HTML, etc. Tu respuesta 
    no captura este principio fundamental de diseño de software.

\end{enumerate}

\section*{Resumen General}

En general demuestras una comprensión correcta de los conceptos de DNS evaluados (MX y su prioridad). 
Sin embargo, en la pregunta sobre diseño de software, la respuesta carece de la profundidad técnica 
esperada, aunque la intuición general no es incorrecta. Se recomienda reforzar los principios de 
separación de responsabilidades (Separation of Concerns) en el diseño de funciones. \\
\textbf{CALIFICACIóN: 8.33 en base a 10 puntos.}

\end{document}